\documentclass{unigrazpub}
\usepackage{microtype}
\usepackage{booktabs}
\usepackage{hyperref}
\hypersetup{colorlinks=true,linkcolor=black,urlcolor=blue}
\title{A Practical Guide to Open Field Notes}
\author{Your Name}
\date{2026}
\begin{document}
\maketitle
\chapter*{Vorwort}
Dieses kurze Beispiel zeigt den Aufbau einer Monografie mit der Klasse
\texttt{unigrazpub}. Der Inhalt ist neu formuliert und kann durch ein eigenes
Manuskript ersetzt werden. Die Struktur eignet sich für ein Buch oder einen
herausgegebenen Band der Universitätsbibliothek.
\tableofcontents
\chapter{Ausgangslage}
Offene wissenschaftliche Notizen brauchen eine nachvollziehbare Ordnung. Ein
Kapitel kann die Fragestellung, die Methode und die Grenzen einer Untersuchung
gemeinsam dokumentieren.
\section{Arbeitsablauf}
\begin{enumerate}
  \item Frage und Zielgruppe festlegen.
  \item Quellen und Beobachtungen mit Datum erfassen.
  \item Ergebnisse, Unsicherheiten und nächste Schritte trennen.
\end{enumerate}
\chapter{Ergebnisse}
Die folgende Tabelle zeigt eine kleine, fiktive Auswertung. Ersetzen Sie sie
durch eigene Daten und ergänzen Sie die Quellenverwaltung des Projekts, sobald
das Manuskript wächst.
\begin{table}[ht]
\centering
\caption{Beispielhafte Arbeitsstände}
\begin{tabular}{lrr}
\toprule
Abschnitt & Entwürfe & geprüft \\
\midrule
Methode & 3 & 2 \\
Ergebnisse & 5 & 4 \\
Anhang & 2 & 1 \\
\bottomrule
\end{tabular}
\end{table}
\chapter{Ausblick}
Ein gutes Buch macht nicht nur die fertige Aussage sichtbar, sondern auch den
Weg, auf dem sie entstanden ist. Nutzen Sie den Ausblick für offene Fragen,
weitere Literatur und Hinweise zur nächsten Ausgabe.
\chapter*{Impressum}
Verlag, Ausgabe, Satz und Bildnachweise können hier ergänzt werden. Die
Bibliothek empfiehlt, die Metadaten und die Lizenz jedes hinzugefügten Assets
gemeinsam mit dem Manuskript zu archivieren.
\end{document}
